% Generated draft: requires mathematical and source review.
\documentclass[11pt]{article}
\usepackage[margin=25mm]{geometry}
\usepackage{fontspec}
\setmainfont{Latin Modern Roman}
\usepackage{amsmath,amssymb,mathtools,unicode-math}
\setmathfont{Latin Modern Math}
\usepackage{enumitem,needspace,xurl,hyperref}
\setlength{\parindent}{0pt}
\setlength{\parskip}{5pt}
\setlength{\emergencystretch}{3em}
\newcommand{\N}{\mathbb{N}}
\newcommand{\Z}{\mathbb{Z}}
\newcommand{\Q}{\mathbb{Q}}
\newcommand{\R}{\mathbb{R}}
\newcommand{\C}{\mathbb{C}}
\newcommand{\F}{\mathbb{F}}
\newcommand{\A}{\mathbb{A}}
\newcommand{\PP}{\mathbb{P}}
\newcommand{\E}{\mathbb{E}}
\newcommand{\eps}{\varepsilon}
\newcommand{\dd}{\,\mathrm{d}}
\begin{document}
\section*{UnsolvedMath ID: 800947}
\textbf{Source problem code:} catalogue entry
\subsection{Definitions and mathematical statement}
Let \((M,g)\) be a complete, simply connected Riemannian manifold, and let \(d_g\) denote its Riemannian distance. The condition that \(M\) has no conjugate points means that no nonzero Jacobi field along a geodesic segment vanishes at both endpoints of that segment. A parametrized geodesic \(\gamma\) satisfies \(\nabla_{\dot\gamma}\dot\gamma=0\).

The supplied question is the following. For every such manifold without conjugate points and every pair of geodesics \(\alpha,\beta\), whose parameters are defined at \(0\) and for all positive times, does
\[
\alpha(0)=\beta(0)
\quad\Longrightarrow\quad
\lim_{t\to+\infty}d_g\bigl(\alpha(t),\beta(t)\bigr)=+\infty
\]
hold? Here the limit means that for every \(R>0\) there exists \(T\geq0\) such that
\[
 t\geq T\quad\Longrightarrow\quad d_g\bigl(\alpha(t),\beta(t)\bigr)>R.
\]

\textbf{Unresolved qualifications.} The supplied statement does not require \(\alpha\) and \(\beta\) to be distinct, does not exclude constant geodesics, and does not prescribe unit speed or another speed normalization. In particular, its stated hypotheses allow coincident parametrized geodesics. These omissions prevent an unambiguous formulation of the intended divergence question. No distinctness or normalization hypothesis has been inserted into the displayed source formulation.
\subsection{Short English statement}
The source asks whether two geodesics starting at the same point in a complete, simply connected Riemannian manifold without conjugate points must become arbitrarily far apart as time increases. The intended restrictions on distinctness, nonconstancy, and parametrization are unspecified and require clarification.
\subsection{Sources}
\begin{itemize}
\item Supplied local UnsolvedMath registry record, permanent ID \(800947\), exact source problem code \texttt{catalogue entry}; title: \emph{Open Problems and Questions about Geodesics --- Question 27}. The record attributes the problem to Burns and Matveev and gives the proposed year as 2019.
\item Burns and Matveev, \emph{Open Problems and Questions about Geodesics} (2019, as identified in the supplied record), Question 27. Supplied source URL: \url{https://arxiv.org/abs/1308.5417}. The provenance records access on 2026-10-16 and extraction by \texttt{arxiv\_latex}; these are supplied provenance claims, not independently verified access or extraction.
\item The registry labels the problem \emph{open}. This is a supplied status label, not an independently verified assessment of the literature or mathematical status.
\end{itemize}
\end{document}
